% يجب الترجمة بمحرك LuaLaTeX  (Menu -> Compiler -> LuaLaTeX)
\documentclass[12pt, a4paper]{article}
\usepackage[a4paper, top=1.2cm, bottom=1.4cm, left=1.4cm, right=1.4cm]{geometry}
\usepackage{amsmath, amssymb}
\usepackage{fontspec}
\usepackage[bidi=basic]{babel}
\babelprovide[main, import, onchar=ids fonts]{arabic}
\babelprovide[import, onchar=ids fonts]{english}
\babelfont{rm}{Noto Serif}
\babelfont[arabic]{rm}{Amiri}
\usepackage{enumitem}
\usepackage{tabularx}
\usepackage{array}
\usepackage{graphicx}
\usepackage{tikz}
\usepackage{eso-pic}

\setlength{\parindent}{0pt}
\newcommand{\LTRfig}[1]{\mbox{\foreignlanguage{english}{#1}}}

% عنوان سؤال داخل شريط رمادي مع العلامة على اليسار
\newcommand{\question}[2]{%
  \par\vspace{0.2cm}\noindent
  \colorbox{gray!15}{\parbox{\dimexpr\linewidth-2\fboxsep\relax}{\textbf{\large #1}\hfill\mbox{\textbf{(#2)}}}}%
  \par\vspace{0.12cm}}
% علامة الفرع في نهاية السطر
\newcommand{\pts}[1]{\hfill\mbox{\textbf{(#1)}}}
% سطر إجابة قصير
\newcommand{\blank}[1][5cm]{\rule{#1}{0.4pt}}
% مربع فارغ للكتابة فيه
\newcommand{\sqbox}{\fbox{\rule{0pt}{0.45cm}\hspace{0.8cm}}}

% ---- شبكة العمليات العمودية (الأرقام من اليسار إلى اليمين كما في الكتاب) ----
% \vgrid{عدد الأعمدة}{عدد الصفوف}
\newcommand{\vgrid}[2]{%
  \draw[gray!70] (0,0) grid[step=0.75] (#1*0.75,-#2*0.75);
  \draw[very thick] (0,-1.5) -- (#1*0.75,-1.5);}
% \digits{الصف}{أول عمود}{الأرقام مفصولة بفواصل}
\newcommand{\digits}[3]{%
  \foreach \d [count=\i from 0] in {#3}
    \node[font=\large] at ({(#2+\i)*0.75+0.375},{-(#1-1)*0.75-0.375}) {\d};}

\begin{document}

% إطار الصفحات
\AddToShipoutPictureBG{%
  \LTRfig{%
  \begin{tikzpicture}[remember picture, overlay]
    \draw[line width=1.2pt] ([shift={(0.7cm,-0.7cm)}]current page.north west) rectangle ([shift={(-0.7cm,0.7cm)}]current page.south east);
  \end{tikzpicture}}%
}

% ============================ الترويسة ============================
\noindent
\begin{minipage}[t]{0.37\textwidth}
\vspace{0pt}
\textbf{المملكة الأردنية الهاشمية} \\
\textbf{وزارة التربية والتعليم} \\
\textbf{مدرسة عنبة الثانوية الشاملة للبنات} \\
\textbf{العام الدراسي: 2025 / 2026}
\end{minipage}%
\begin{minipage}[t]{0.26\textwidth}
\vspace{0pt}
\centering
\IfFileExists{logo.png}{\includegraphics[width=2.2cm]{logo.png}}{}
\end{minipage}%
\begin{minipage}[t]{0.37\textwidth}
\vspace{0pt}
\textbf{المبحث: الرياضيات} \\
\textbf{الصف: الخامس} \\
\textbf{الزمن: ساعتان} \\
\textbf{العلامة الكلية: 33 علامة}
\end{minipage}

\vspace{0.2cm}
\hrule height 1.2pt
\vspace{0.06cm}
\hrule height 0.4pt
\vspace{0.15cm}

\begin{center}
\textbf{\LARGE اختبار نهاية الوحدة}
\end{center}
\vspace{-0.15cm}

% بيانات الطالبة
\noindent\fbox{\parbox{\dimexpr\linewidth-2\fboxsep-2\fboxrule\relax}{%
\vspace{0.15cm}
\textbf{الاسم:} \rule{6.2cm}{0.4pt} \hfill
\textbf{الشعبة:} \rule{2cm}{0.4pt} \hfill
\textbf{التاريخ:} \rule{0.9cm}{0.4pt} / \rule{0.9cm}{0.4pt} / 2026
\vspace{0.1cm}}}

\vspace{0.15cm}
% جدول توزيع العلامات
\begin{center}
\renewcommand{\arraystretch}{1.2}
\begin{tabular}{|>{\bfseries}c|*{7}{>{\centering\arraybackslash}p{1.1cm}|}>{\centering\arraybackslash}p{1.4cm}|}
\hline
السؤال & الأول & الثاني & الثالث & الرابع & الخامس & السادس & السابع & المجموع \\ \hline
العلامة & 7 & 5 & 6 & 4 & 6 & 3 & 2 & \textbf{33} \\ \hline
علامة الطالبة & & & & & & & & \\ \hline
\end{tabular}
\end{center}

\vspace{-0.15cm}
\noindent\textbf{تعليمات:} أجب عن جميع الأسئلة، واكتب خطوات الحل في المكان المخصص. الآلة الحاسبة غير مسموحة.

% ============================ السؤال الأول ============================
\question{السؤال الأول: المفاهيم والأعداد}{7 علامات}

\textbf{أ)} اكتب القيمة المنزلية للرقم الذي تحته خط في العدد: \quad $7\underline{8}3596015$
\quad \blank[4.5cm] \pts{علامتان}
\vspace{0.3cm}

\textbf{ب)} اكتب العدد (مليون ومئة وأربعة وسبعون ألفاً وأربعمئة وثلاثة وتسعون) بالصيغتين القياسية والتحليلية. \pts{علامتان}
\par\vspace{0.35cm}
الصيغة القياسية: \blank[4cm] \par\vspace{0.4cm}
الصيغة التحليلية: \blank[13cm]
\vspace{0.3cm}

\textbf{ج)} تبلغ المسافة بين الأرض والشمس $149598428$ كم تقريباً. اكتب هذا العدد بالصيغة اللفظية. \pts{علامتان}
\par\vspace{0.6cm}
\rule{\linewidth}{0.4pt}
\par\vspace{0.6cm}
\rule{\linewidth}{0.4pt}
\vspace{0.3cm}

\textbf{د)} كتبت سميرة العدد $35003936$ بالصيغة التحليلية:\quad $6 + 30 + 900 + 3000 + 50000 + 30000000$.
اكتشف خطأ سميرة وصححه. \pts{علامة واحدة}
\par\vspace{0.5cm}
الخطأ: \blank[6cm] \hfill التصحيح: \blank[7cm]

% ============================ السؤال الثاني ============================
\question{السؤال الثاني: المقارنة والترتيب}{5 علامات}

\textbf{أ)} ضع الرمز المناسب ($<$ ، $>$ ، $=$) في المربع لتصبح العبارة صحيحة: \pts{علامتان}
\par\vspace{0.15cm}
\begin{tabularx}{\linewidth}{XX}
\textbf{1)} \ $3759120 \ \ \sqbox \ \ 3456210$ & \textbf{2)} \ $4589000 \ \ \sqbox \ \ 4589000$
\end{tabularx}
\vspace{0.2cm}

\textbf{ب)} اكتب رقماً مناسباً في المربع لتصبح الجملة صحيحة: \quad
$9124382 < 91243\,\sqbox\,2$ \pts{علامة واحدة}
\vspace{0.3cm}

\textbf{ج)} رتّب الأعداد الآتية تصاعدياً: \quad $53241290$ \ ، \ $53641025$ \ ، \ $53279243$ \pts{علامتان}
\par\vspace{0.5cm}
\blank[4cm] \ \ ، \ \ \blank[4cm] \ \ ، \ \ \blank[4cm]

\newpage
% ============================ السؤال الثالث ============================
\question{السؤال الثالث: الجمع والطرح}{6 علامات}

\noindent
\begin{minipage}[t]{0.48\linewidth}
\textbf{أ)} أجد ناتج الجمع في المكان المخصص: \pts{3 علامات}
\par\vspace{0.2cm}\centering
\LTRfig{%
\begin{tikzpicture}
  \vgrid{8}{4}
  \digits{1}{1}{4,6,6,1,2,4,3}
  \digits{2}{1}{2,3,6,7,9,8,1}
  \node[font=\Large\bfseries] at (-0.4,-1.125) {$+$};
\end{tikzpicture}}
\end{minipage}\hfill\vrule\hfill
\begin{minipage}[t]{0.48\linewidth}
\textbf{ب)} أجد ناتج الطرح في المكان المخصص: \pts{3 علامات}
\par\vspace{0.2cm}\centering
\LTRfig{%
\begin{tikzpicture}
  \vgrid{8}{4}
  \digits{1}{1}{9,7,1,8,3,9,6}
  \digits{2}{1}{8,7,4,5,6,1,3}
  \node[font=\Large\bfseries] at (-0.4,-1.125) {$-$};
\end{tikzpicture}}
\end{minipage}

% ============================ السؤال الرابع ============================
\question{السؤال الرابع: الضرب}{4 علامات}

\noindent
\begin{minipage}[t]{0.5\linewidth}
أجد ناتج الضرب عمودياً في المكان المخصص:
\end{minipage}\hfill
\begin{minipage}[t]{0.45\linewidth}
\vspace{0pt}\centering
\LTRfig{%
\begin{tikzpicture}
  \vgrid{6}{4}
  \digits{1}{2}{3,6,2,1}
  \digits{2}{5}{8}
  \node[font=\Large\bfseries] at (-0.4,-1.125) {$\times$};
\end{tikzpicture}}
\end{minipage}

% ============================ السؤال الخامس ============================
\question{السؤال الخامس: المسائل الكلامية والقسمة}{6 علامات}

\noindent
\begin{minipage}[t]{0.48\linewidth}
\textbf{أ)} ينام دب الكوالا $18$ ساعة يومياً. فكم ساعة ينام في العام الواحد (اعتبر العام $365$ يوماً)؟ \pts{علامتان}
\par\vspace{2.5cm}
الجواب: \blank[4.5cm]
\end{minipage}\hfill\vrule\hfill
\begin{minipage}[t]{0.48\linewidth}
\textbf{ب)} وزّع رجل $986$ ديناراً زكاة على $17$ فقيراً بالتساوي. كم نصيب كل منهم؟ \pts{علامتان}
\par\vspace{2.5cm}
الجواب: \blank[4.5cm]
\end{minipage}

\vspace{0.35cm}
\textbf{ج)} أراد مدير مدرسة نقل $445$ طالباً في حافلات سعتها $35$ راكباً. كم حافلة يحتاج؟ وماذا يمثل الباقي؟ (استخدم القسمة الطويلة) \pts{علامتان}
\par\vspace{2.4cm}
كم حافلة يحتاج؟ \blank[4cm] \hfill وماذا يمثل الباقي؟ \blank[5.5cm]

% ============================ السؤالان السادس والسابع متجاوران ============================
\vspace{0.2cm}
\noindent
\begin{minipage}[t]{0.48\linewidth}
\question{السؤال السادس: الأعداد الصحيحة}{3 علامات}
يتكون برج في عمّان من $50$ طابقاً فوق الأرض، و$7$ طوابق تحتها. ما العدد المناسب لوصف الطابق السابع تحت الأرض؟
\par\vspace{1.4cm}
الجواب: \blank[4.5cm]
\end{minipage}\hfill\vrule\hfill
\begin{minipage}[t]{0.48\linewidth}
\question{السؤال السابع: قابلية القسمة}{علامتان}
\textbf{أ)} ابحث قابلية قسمة العدد $5124$ على $4$، مع التبرير. \pts{علامة واحدة}
\par\vspace{1.3cm}
\textbf{ب)} التقط مصور $144$ صورة. هل يمكنه وضعها في $6$ لوحات جدارية بالتساوي؟ \pts{علامة واحدة}
\par\vspace{1.1cm}
\end{minipage}

\vfill
\begin{center}
\vspace{-0.3cm}
\textbf{\large انتهت الأسئلة، مع تمنياتي لكم بالتوفيق}
\end{center}

\end{document}
